% =============================== PORTADA ===============================
\thispagestyle{empty}
\setcounter{page}{0}
\vspace*{-1.2cm}
\noindent\begin{tikzpicture}[remember picture,overlay]
  \shade[left color=navy,right color=regla!60!white] (current page.north west) rectangle ([yshift=-3.2mm]current page.north east);
  \shade[left color=navy,right color=regla!60!white] (current page.south west) rectangle ([yshift=3.2mm]current page.south east);
\end{tikzpicture}%
\begin{tikzpicture}
  \useasboundingbox (0,-1.15) rectangle (\textwidth,0.75);
  \node[anchor=west,inner sep=0pt] at (0,0) {\includegraphics[width=0.22\textwidth]{logo-usal.png}};
  \node[anchor=east,inner sep=0pt] at (\textwidth,0) {\includegraphics[width=0.15\textwidth]{logo-ciencias.png}};
  \node[align=center,inner sep=0pt] at (0.535\textwidth,0)
    {{\color{navy}\large\espaciado[150]{Universidad de Salamanca}}\\[4pt]
     {\color{acento}\espaciado[110]{Facultad de Ciencias}\enspace\textperiodcentered\enspace\espaciado[110]{Grado en Física}}};
  \draw[regla,line width=0.6pt] (0,-1.1) -- (\textwidth,-1.1);
\end{tikzpicture}

\vspace{0.8cm}

\noindent{\color{acento}\large\espaciado[150]{Termodinámica I}\enspace\textperiodcentered\enspace\espaciado[150]{Tema 1}}\\[8pt]
{\color{navy}\fontsize{30}{36}\selectfont\bfseries Conceptos básicos}\\[4pt]
\begin{tikzpicture}
  \shade[left color=navy,right color=white] (0,0) rectangle (0.62\textwidth,1.6mm);
\end{tikzpicture}

\vspace{0.5cm}
\noindent{\large\itshape Sistemas y paredes, equilibrio y procesos, ley cero y temperatura empírica, ecuaciones de estado y coeficientes de respuesta $\alpha$, $\kappa_T$ y $\beta$}

\vspace{0.7cm}
\noindent
\begin{minipage}[t]{0.5\textwidth}
  {\color{acento}\espaciado[180]{Autor}}\\
  {\large Daniel Martín Guerrero\,\textsuperscript{\color{acento}a}}
\end{minipage}%
\begin{minipage}[t]{0.5\textwidth}
  {\color{acento}\espaciado[180]{Profesor}}\\
  {\large Julián González Ayala\,\textsuperscript{\color{acento}a}}
\end{minipage}

\vspace{4pt}
\noindent{\color{gris}\textsuperscript{a}\,Grado en Física, Facultad de Ciencias, Universidad de Salamanca, Plaza de la Merced s/n, 37008 Salamanca, España}

\vspace{0.5cm}
\begin{tcolorbox}[cajabase,top=10pt,bottom=10pt]
\begin{minipage}[t]{0.52\linewidth}
  {\color{acento}\espaciado[200]{Resumen}}\par\smallskip
  \small
  Establecemos el lenguaje de la termodinámica de equilibrio: sistema, frontera y paredes,
  variables extensivas e intensivas y la noción de estado. Distinguimos equilibrio de estado
  estacionario y clasificamos los procesos (cuasiestático, reversible, irreversible). La ley cero
  convierte el equilibrio térmico en una relación de equivalencia y garantiza la existencia de una
  temperatura empírica. Por último, la ecuación de estado $f(p,V,T)=0$ y sus derivadas, los
  coeficientes $\alpha$, $\kappa_T$ y $\beta$, ligados por $\alpha=p\beta\kappa_T$, permiten
  reconstruir la ecuación de estado integrando una diferencial exacta. Los seminarios resuelven,
  paso a paso y con todo el desarrollo simbólico, los dieciséis problemas del tema y la cuestión 1
  del parcial.
\end{minipage}\hfill
\begin{minipage}[t]{0.44\linewidth}\centering
  \vspace{2pt}
  \begin{tikzpicture}
  \begin{axis}[view={30}{25},width=6.9cm,height=6.6cm,
     xmin=0.6,xmax=4.2,ymin=1,ymax=4,zmin=0,zmax=7,
     xtick=\empty,ytick=\empty,ztick=\empty,
     axis lines=left,axis line style={navy,-{Stealth[length=5pt]}},
     xlabel={$V$},ylabel={$T$},zlabel={$p$},ylabel style={anchor=north west},
     colormap={estado}{rgb255(0cm)=(232,239,248); rgb255(1cm)=(157,183,210); rgb255(2cm)=(58,110,165); rgb255(3cm)=(11,37,69)},
     point meta min=0.3,point meta max=5.6,clip=false]
    \addplot3[surf,shader=interp,domain=0.6:4.2,y domain=1:4,samples=34,samples y=22] {min(y/x,6)};
    \addplot3[mesh,draw=white,draw opacity=0.3,line width=0.3pt,domain=0.6:4.2,y domain=1:4,samples=13,samples y=7] {min(y/x,6)};
    \pgfplotsinvokeforeach{1.5,2.5,3.5}{
      \addplot3[white,line width=1.5pt,domain=#1/6:4.2,samples=50,samples y=1] (x,#1,{#1/x});
    }
    \addplot3[rojo,line width=1.4pt,domain=0.75:1.75,samples=30,samples y=1] (x,{2*x},2);
    \fill[rojo] (axis cs:0.75,1.5,2) circle (1.6pt);
    \fill[rojo] (axis cs:1.75,3.5,2) circle (1.6pt);
    \coordinate (iso) at (axis cs:3.2,3.5,{3.5/3.2});
    \coordinate (isob) at (axis cs:1.5,3.0,2);
  \end{axis}
  \draw[gris,line width=0.4pt] (iso) -- ++(0.7,0.65) node[anchor=south west,inner sep=1pt,font=\scriptsize,navy] {isotermas};
  \draw[gris,line width=0.4pt] (isob) -- ++(0.9,1.35) node[anchor=south west,inner sep=1pt,font=\scriptsize,rojo] {isobara};
  \end{tikzpicture}\par
  {\scriptsize\itshape\color{gris} Superficie de estado $pV=nRT$ del gas ideal: cada punto es un estado de
  equilibrio; sus cortes a $T$ constante son las isotermas y a $p$ constante, las isobaras.\par}
\end{minipage}
\end{tcolorbox}

\noindent{\color{acento}\espaciado[200]{Palabras clave}}\enspace
sistema termodinámico\sepdot paredes y ligaduras\sepdot extensiva e intensiva\sepdot equilibrio\sepdot
proceso cuasiestático\sepdot ley cero\sepdot temperatura empírica\sepdot ecuación de estado\sepdot
coeficientes termomecánicos\sepdot relación cíclica\par
\smallskip
\noindent{\color{acento}\espaciado[200]{Clasificación PACS}}\enspace
05.70.$-$a (Termodinámica), 05.70.Ce (Funciones termodinámicas y ecuaciones de estado), 65.40.De (Dilatación térmica)

\vspace{0.4cm}
{\color{regla}\hrule height 0.6pt}
\vspace{2pt}
\[
  pV=nRT\qquad
  \alpha=\frac1V\pd{V}{T}{p}\qquad
  \kappa_T=-\frac1V\pd{V}{p}{T}\qquad
  \alpha=p\,\beta\,\kappa_T\qquad
  \dd\ln V=\alpha\,\dd T-\kappa_T\,\dd p
\]
\vspace{-2pt}
{\color{regla}\hrule height 0.6pt}

\vfill
\begin{center}
  {\color{navy}\large\espaciado[220]{Apuntes formalizados con seminarios prácticos}}\\[4pt]
  {\color{acento}\espaciado[200]{Salamanca, curso 2026/2027}}
\end{center}
\vspace*{0.1cm}
